% ============================================================================
% Rapport Technique — Drone de surface à voile sur T-Beam ESP32
% Langue : Français (traduction de la version espagnole)
% Auteur (code original) : Mohamed EL JILY
% Auteur (refonte / firmware actuel) : Facundo Arito
% Compilation : pdflatex -interaction=nonstopmode rapport_projet_fr.tex (x2)
% ============================================================================

\documentclass[12pt,a4paper]{article}

% ── Encodage (pdflatex) ──────────────────────────────────────────────────────
\usepackage[utf8]{inputenc}
\usepackage[T1]{fontenc}
\usepackage{lmodern}
\usepackage{textcomp}

% ── Langue (sans babel-french pour ne pas dépendre de texlive-lang) ──────────
\renewcommand{\contentsname}{Table des mati\`eres}
\renewcommand{\abstractname}{R\'esum\'e}
\renewcommand{\appendixname}{Annexe}
\renewcommand{\refname}{R\'ef\'erences}
\renewcommand{\figurename}{Figure}
\renewcommand{\tablename}{Tableau}

% ── Mise en page ─────────────────────────────────────────────────────────────
\usepackage[top=2.5cm, bottom=2.5cm, left=2.8cm, right=2.8cm]{geometry}
\usepackage{parskip}
\usepackage{microtype}
\usepackage{setspace}
\onehalfspacing

% ── Couleurs ─────────────────────────────────────────────────────────────────
\usepackage[dvipsnames]{xcolor}
\definecolor{codebg}{HTML}{1E1E2E}
\definecolor{codefg}{HTML}{CDD6F4}
\definecolor{commentfg}{HTML}{6C7086}
\definecolor{keywordfg}{HTML}{CBA6F7}
\definecolor{stringfg}{HTML}{A6E3A1}
\definecolor{numberfg}{HTML}{FAB387}
\definecolor{warnbg}{HTML}{FFF8E1}
\definecolor{warnborder}{HTML}{F59E0B}
\definecolor{okbg}{HTML}{F0FDF4}
\definecolor{okborder}{HTML}{16A34A}
\definecolor{infobg}{HTML}{EFF6FF}
\definecolor{infoborder}{HTML}{2563EB}
\definecolor{tableheadbg}{HTML}{374151}
\definecolor{tableheadfg}{HTML}{FFFFFF}
\definecolor{tablerowalt}{HTML}{F9FAFB}

% ── Boîtes d'avis ────────────────────────────────────────────────────────────
\usepackage{tcolorbox}
\tcbuselibrary{skins, breakable}

\newtcolorbox{warningbox}{
  colback=warnbg, colframe=warnborder,
  fonttitle=\bfseries, title={\color{warnborder}$\blacktriangle$\ Probl\`eme},
  breakable, boxrule=1pt, left=6pt, right=6pt
}
\newtcolorbox{okbox}{
  colback=okbg, colframe=okborder,
  fonttitle=\bfseries, title={\color{okborder}$\checkmark$\ Solution},
  breakable, boxrule=1pt, left=6pt, right=6pt
}
\newtcolorbox{infobox}[1][]{
  colback=infobg, colframe=infoborder,
  fonttitle=\bfseries, title={\color{infoborder}$\blacksquare$\ #1},
  breakable, boxrule=1pt, left=6pt, right=6pt
}

% ── Listings (code source) ───────────────────────────────────────────────────
\usepackage{listings}
\usepackage{lstautogobble}

\lstdefinestyle{cpp}{
  language=C++,
  backgroundcolor=\color{codebg},
  basicstyle=\ttfamily\footnotesize\color{codefg},
  keywordstyle=\color{keywordfg}\bfseries,
  stringstyle=\color{stringfg},
  commentstyle=\color{commentfg}\itshape,
  numberstyle=\tiny\color{numberfg},
  numbers=left, numbersep=8pt,
  breaklines=true, breakatwhitespace=true,
  frame=single, rulecolor=\color{codebg!40!black},
  tabsize=4, showstringspaces=false,
  autogobble=true,
  morekeywords={uint8_t,uint16_t,uint32_t,int8_t,int16_t,int32_t,bool,
                constexpr,noexcept,nullptr,override,IRAM_ATTR,portMUX_TYPE},
}
\lstdefinestyle{shell}{
  language=bash,
  backgroundcolor=\color{codebg},
  basicstyle=\ttfamily\footnotesize\color{codefg},
  breaklines=true, frame=single, rulecolor=\color{codebg!40!black},
  showstringspaces=false,
}
\lstdefinestyle{json}{
  backgroundcolor=\color{codebg},
  basicstyle=\ttfamily\footnotesize\color{codefg},
  breaklines=true, frame=single, rulecolor=\color{codebg!40!black},
  showstringspaces=false,
}

% ── Tableaux ─────────────────────────────────────────────────────────────────
\usepackage{booktabs}
\usepackage{array}
\usepackage{longtable}
\usepackage{colortbl}
\usepackage{multirow}
\newcolumntype{L}[1]{>{\raggedright\arraybackslash}p{#1}}
\newcolumntype{C}[1]{>{\centering\arraybackslash}p{#1}}

% ── Hyperliens ───────────────────────────────────────────────────────────────
\usepackage[hidelinks]{hyperref}
\hypersetup{
  pdftitle={Rapport Technique -- Drone de surface à voile},
  pdfauthor={Facundo Arito},
  bookmarksnumbered=true,
}

% ── Mise en forme des titres ─────────────────────────────────────────────────
\usepackage{titlesec}
\titleformat{\section}{\Large\bfseries\color{tableheadbg}}{\thesection.}{0.7em}{}[\vspace{-0.3em}\color{infoborder}\rule{\linewidth}{1.5pt}]
\titleformat{\subsection}{\large\bfseries\color{tableheadbg!80!black}}{\thesubsection.}{0.7em}{}
\titleformat{\subsubsection}{\normalsize\bfseries}{\thesubsubsection.}{0.6em}{}

% ── Listes ───────────────────────────────────────────────────────────────────
\usepackage{enumitem}
\setlist{itemsep=2pt, topsep=4pt}

% ── En-têtes et pieds de page ────────────────────────────────────────────────
\usepackage{fancyhdr}
\setlength{\headheight}{15pt}
\pagestyle{fancy}
\fancyhf{}
\fancyhead[L]{\small Drone de surface à voile --- Rapport technique}
\fancyhead[R]{\small \today}
\fancyfoot[C]{\small\thepage}
\renewcommand{\headrulewidth}{0.4pt}

% ── Utilitaires ──────────────────────────────────────────────────────────────
\usepackage{graphicx}
\usepackage{amsmath}
\usepackage{amssymb}

\newcommand{\gpio}[1]{\texttt{GPIO#1}}
\newcommand{\us}[1]{\texttt{#1\,\textmu s}}
\newcommand{\file}[1]{\texttt{#1}}
\newcommand{\done}{\textcolor{okborder}{\textbf{[OK]}}}
\newcommand{\todo}{\textcolor{warnborder}{\textbf{[\ \ ]}}}

% ─────────────────────────────────────────────────────────────────────────────
\begin{document}

% ── Page de titre ────────────────────────────────────────────────────────────
\begin{titlepage}
  \centering
  \vspace*{2cm}
  {\Huge\bfseries Drone de surface à voile\par}
  \vspace{0.5cm}
  {\Large Rapport technique --- D\'eveloppement du firmware et de l'\'electronique\par}
  \vspace{2cm}
  \rule{\linewidth}{1pt}
  \vspace{0.5cm}

  \begin{tabular}{ll}
    \textbf{Plateforme :} & LilyGO T-Beam V1.1 (ESP32 + SX1276 + GPS u-blox + AXP192) \\[4pt]
    \textbf{Langage :}    & Arduino / C++, cible ESP32 Arduino Core 3.x \\[4pt]
    \textbf{Membres du groupe : } & Facundo Arito \\[4pt]
    \textbf{Date :}       & \today \\
  \end{tabular}

  \vspace{0.5cm}
  \rule{\linewidth}{1pt}
  \vspace{2cm}

  \begin{abstract}
    Ce document retrace l'int\'egralit\'e du d\'eveloppement informatique et
    \'electronique du drone de surface à voile bas\'e sur le module T-Beam
    d'Espressif/LilyGO. Il couvre le concept physique du v\'ehicule, le code
    original et ses limitations, la r\'e\'ecriture compl\`ete de l'architecture
    logicielle, les deux r\'evisions du PCB, le mode manuel unifi\'e avec
    propulsion bidirectionnelle, la navigation autonome à voile, l'interface de
    station sol (IHM) et, de mani\`ere d\'etaill\'ee, \textbf{tous les probl\`emes}
    rencontr\'es au cours du projet ainsi que les solutions mises en place. Une
    section finale r\'eunit les probl\`emes encore ouverts et les options
    d'am\'elioration propos\'ees pour les r\'esoudre.
  \end{abstract}

  \vfill
  \textit{Projet de recherche (PER) --- Drone autonome d'acquisition de donn\'ees marines}
\end{titlepage}

\tableofcontents
\newpage

% =============================================================================
\section{Guide d'utilisation pour le prochain groupe de travail}
\label{sec:guia}
% =============================================================================

Cette section est un \textbf{guide pratique de mise en marche} destin\'e au groupe
suivant qui reprendra le projet. Elle explique quel mat\'eriel est n\'ecessaire,
comment allumer et piloter le drone, comment programmer et connecter les deux
cartes T-Beam, et comment lancer l'interface de station sol (IHM) aussi bien sous
Linux que sous Windows. Il est recommand\'e de la lire \textbf{en entier} avant le
premier allumage.

\subsection{Mat\'eriel n\'ecessaire}

\begin{longtable}{L{4.5cm} L{8.5cm}}
  \toprule
  \rowcolor{tableheadbg}
  \textcolor{tableheadfg}{\textbf{\'El\'ement}} &
  \textcolor{tableheadfg}{\textbf{D\'etail}} \\
  \midrule
  \endhead
  \bottomrule
  \endfoot
  Drone (``bateau'') & T-Beam mont\'ee dans la coque sur le PCB V2. Convertisseur
                       USB-s\'erie CP2104, num\'ero de s\'erie \texttt{02126CF1} \\[3pt]
  Transcepteur (station sol) & Deuxi\`eme T-Beam avec le sketch
                       \file{transceiver.ino}. CP2104 s\'erie \texttt{01C00B54} \\[3pt]
  Radio RC & \'Emetteur Pro-Tronik PTR-6A + r\'ecepteur R8X (2,4~GHz) \\[3pt]
  Batterie & LiPo 2S \textbf{charg\'ee} ($\approx$8,4~V) + porte-batteries de la
             coque. \textbf{Indispensable} pour un fonctionnement stable \\[3pt]
  Ordinateur & Linux (recommand\'e) ou Windows, avec Docker, Python 3 et navigateur \\[3pt]
  C\^ables USB & Un de bonne qualit\'e pour le transcepteur~; un autre pour
                programmer / alimenter le bateau \\
\end{longtable}

\subsection{Mise en marche du drone (bateau)}

\begin{enumerate}
  \item \textbf{Charger la LiPo} compl\`etement et la connecter au \textbf{bornier du
    PCB} du bateau (le GND est indiqu\'e sur le PCB).
  \item \textbf{Allumer d'abord l'\'emetteur RC} et v\'erifier que le s\'electeur de
    mode (CH5) est en \textbf{Manuel} (position basse). Le \textbf{premier essai doit
    toujours \^etre manuel} afin de tester tous les actionneurs.
  \item \textbf{Mettre le bateau sous tension.} Connecter l'USB de la T-Beam (soit à
    l'alimentation du PCB, soit à un PC pour lire le port s\'erie). Au d\'emarrage, il
    imprime sur le port s\'erie des lignes d'\'etat (p.\,ex.\ \texttt{[LORA] OK}).
  \item \textbf{Attendre le fix GPS :} la LED \texttt{TIMEPULSE} du NEO-6M clignote à
    1~Hz quand il y a un fix. Avec la LiPo connect\'ee, le fix arrive en quelques
    secondes (d\'emarrage à chaud)~; sans LiPo, chaque allumage est un d\'emarrage à
    froid d'$\approx$2~min.
  \item \textbf{Piloter en manuel :}
  \begin{itemize}
    \item \textbf{CH5 (mode) :} haut $=$ Automatique, milieu $=$ Sail (inerte, tout
      au neutre), bas $=$ Manuel.
    \item \textbf{CH2 (voile) :} bascule la voile entre b\^abord et tribord ($\pm10^{\circ}$).
    \item \textbf{CH4 (safran) :} positionne le treuil du safran ($\pm90^{\circ}$).
    \item \textbf{CH3 (moteur, bidirectionnel) :} centre $=$ arr\^et, une extr\'emit\'e
      $=$ $+100~\%$ avant, l'autre $=$ $-100~\%$ arri\`ere. N\'ecessite l'ESC en mode
      bidirectionnel (voir point suivant).
  \end{itemize}
\end{enumerate}

\begin{warningbox}
  \textbf{S\'ecurit\'e :} en cas de perte du signal RC, le drone passe en
  \emph{Failsafe}, qui \textbf{navigue de fa\c{c}on autonome}. Comme le mode autonome
  n'a jamais \'et\'e valid\'e sur l'eau, le premier essai doit se faire en manuel, en
  eaux calmes et avec un plan de r\'ecup\'eration (p.\,ex.\ barque ou cordage).
\end{warningbox}

\subsection{Programmation (flash) des T-Beam}

Les deux cartes utilisent le m\^eme noyau (\texttt{FQBN = esp32:esp32:t-beam}). Le
bateau porte le sketch \file{main/}~; le transcepteur porte \file{transceiver/}.

\begin{lstlisting}[style=shell, caption={Compiler et flasher avec arduino-cli}]
# (chemins relatifs au dossier racine du projet = Post-Claude)
# Compiler le firmware du bateau
~/bin/arduino-cli compile --fqbn esp32:esp32:t-beam main

# Flasher le bateau (ajuster le port)
~/bin/arduino-cli upload --port /dev/ttyUSB0 --fqbn esp32:esp32:t-beam main

# Idem pour le transcepteur
~/bin/arduino-cli upload --port /dev/ttyUSB0 --fqbn esp32:esp32:t-beam transceiver
\end{lstlisting}

\begin{itemize}
  \item \textbf{Identifier chaque carte par son num\'ero de s\'erie} (sous Linux :
    \texttt{ls -l /dev/serial/by-id/}). Bateau $=$ \texttt{02126CF1},
    transcepteur $=$ \texttt{01C00B54}.
  \item \textbf{Si le flash \'echoue} (``\emph{chip stopped responding}'', le port se
    r\'e-\'enum\`ere) : c'est un \emph{brownout} par alimentation insuffisante. Entrer
    en mode bootloader manuel (maintenir \textbf{BOOT/IO0} en appuyant sur
    \textbf{RST} ou en reconnectant l'USB). Voir Section~\ref{sec:alim}.
  \item Ne pas avoir de moniteur s\'erie ouvert sur le port pendant la programmation.
\end{itemize}

\subsection{Connexion des deux T-Beam}

\begin{itemize}
  \item Le \textbf{transcepteur} se connecte en \textbf{USB à l'ordinateur} de sol.
    Il appara\^it comme \texttt{/dev/ttyUSBx} sous Linux ou \texttt{COMx} sous Windows.
  \item Le \textbf{bateau ne se connecte pas par c\^able} pendant l'op\'eration : il
    communique avec le transcepteur par \textbf{radio LoRa à 433~MHz}. Les deux
    firmwares utilisent d\'ej\`a les m\^emes param\`etres radio (433~MHz, SF7,
    BW125~kHz).
  \item Le transcepteur retransmet à l'IHM chaque heartbeat du bateau (et lui injecte
    le RSSI mesur\'e au sol) et transmet au bateau les commandes de l'IHM.
\end{itemize}

\begin{warningbox}
  \textbf{Ne pas ouvrir} un moniteur s\'erie (gtkterm, le moniteur de l'IDE Arduino,
  un script pyserial, etc.) sur le port du transcepteur \textbf{pendant que l'IHM
  tourne} : les deux lecteurs se disputent le port et corrompent le JSON de
  t\'el\'em\'etrie.
\end{warningbox}

\subsection{Utilisation de l'IHM sous Linux}

\subsubsection{Programmes à installer (Linux)}

\begin{longtable}{L{4.2cm} L{8.8cm}}
  \toprule
  \rowcolor{tableheadbg}
  \textcolor{tableheadfg}{\textbf{Programme}} &
  \textcolor{tableheadfg}{\textbf{Pour quoi / comment}} \\
  \midrule
  \endhead
  \bottomrule
  \endfoot
  \textbf{Docker Engine} + plugin \texttt{docker compose} & Lance la base de
    donn\'ees \textbf{MongoDB} (image \texttt{mongo:7}, t\'el\'echarg\'ee
    automatiquement la premi\`ere fois). Sous Ubuntu/Debian :
    \texttt{sudo apt install docker.io docker-compose-plugin}. Ajouter l'utilisateur
    au groupe : \texttt{sudo usermod -aG docker \$USER} (puis rouvrir la session)
    pour ne pas utiliser \texttt{sudo} \\[3pt]
  \textbf{Python 3} + \texttt{venv} + \texttt{pip} & Ex\'ecute le serveur et le pont
    s\'erie. \texttt{sudo apt install python3 python3-venv python3-pip} \\[3pt]
  \textbf{Paquets Python} & S'installent dans un environnement virtuel à partir de
    \file{requirements.txt} : \texttt{fastapi}, \texttt{uvicorn},
    \texttt{pymongo[asyncio]}, \texttt{pyserial}, \texttt{python-dotenv} \\[3pt]
  \textbf{Navigateur web} & Pour ouvrir l'interface (Firefox, Chrome, etc.) \\
\end{longtable}

\textbf{Pr\'eparation (une seule fois)} --- cr\'eer l'environnement virtuel et
installer les paquets Python (chemins relatifs à la racine \file{Post-Claude}) :

\begin{lstlisting}[style=shell, caption={Installation initiale de l'IHM sous Linux}]
cd IHM
python3 -m venv .venv
.venv/bin/pip install -r requirements.txt
\end{lstlisting}

\textbf{D\'emarrage et arr\^et} --- connecter le transcepteur en USB et, depuis le
dossier \file{IHM/}, ex\'ecuter :

\begin{lstlisting}[style=shell, caption={Demarrer et arreter l'IHM sous Linux}]
cd IHM
./start_ihm.sh                 # autodetecte le transcepteur par son numero de serie
# ou forcer le port :
./start_ihm.sh /dev/ttyUSB0

# ... a la fin :
./stop_ihm.sh                  # libere le port 5000 et arrete MongoDB (Docker)
\end{lstlisting}

\begin{itemize}
  \item \file{start\_ihm.sh} lance MongoDB (conteneur Docker via
    \texttt{docker compose}), \file{serial\_link.py} (lit/\'ecrit le port s\'erie) et
    \file{webserver.py} (serveur web), et ouvre le navigateur sur
    \texttt{http://localhost:5000}.
  \item Sur la carte : voir la t\'el\'em\'etrie en direct, d\'efinir des waypoints au
    clic, envoyer la route et les commandes (naviguer, arr\^eter, mesurer le vent, etc.).
  \item Commutateur \textbf{Sim / R\'eel} : en ``R\'eel'' il utilise le transcepteur
    physique~; en ``Sim'' il lance un simulateur de bateau (sans mat\'eriel).
  \item Le port du transcepteur est \textbf{autod\'etect\'e} par son num\'ero de
    s\'erie (\texttt{01C00B54}), il n'est donc normalement pas n\'ecessaire de
    l'indiquer.
\end{itemize}

\subsection{Utilisation de l'IHM sous Windows}

L'IHM a \'et\'e d\'evelopp\'ee sous Linux (les scripts de d\'emarrage sont en
\texttt{bash}), c'est pourquoi sous Windows le plus simple est de l'ex\'ecuter dans
WSL2.

\subsubsection{Programmes à installer (Windows)}

\begin{longtable}{L{4.2cm} L{8.8cm}}
  \toprule
  \rowcolor{tableheadbg}
  \textcolor{tableheadfg}{\textbf{Programme}} &
  \textcolor{tableheadfg}{\textbf{Pour quoi / comment}} \\
  \midrule
  \endhead
  \bottomrule
  \endfoot
  \textbf{WSL2} (Ubuntu) & Sous-syst\`eme Linux dans Windows~; permet d'utiliser les
    m\^emes scripts \texttt{bash}. Installer avec \texttt{wsl --install} dans
    PowerShell (en administrateur) \\[3pt]
  \textbf{Docker Desktop} & Lance MongoDB. L'installer et activer
    l'\textbf{int\'egration avec WSL2} dans ses param\`etres \\[3pt]
  \textbf{usbipd-win} & Partage le port USB du transcepteur avec WSL (Windows ne le
    passe pas seul). Installer avec \texttt{winget install usbipd} \\[3pt]
  \textbf{Python 3} + paquets & D\'ej\`a pr\'esent dans l'Ubuntu de WSL~; installer
    les paquets comme sous Linux (\file{requirements.txt}). En Windows natif,
    installer Python depuis \texttt{python.org} \\[3pt]
  \textbf{Navigateur web} & Pour ouvrir \texttt{http://localhost:5000} \\
\end{longtable}

\textbf{Voie recommand\'ee --- WSL2 + Docker Desktop :} dans Ubuntu (WSL) on suit
\textbf{exactement les m\^emes \'etapes que sous Linux} (cr\'eer le \texttt{.venv},
installer \file{requirements.txt}, \file{./start\_ihm.sh}). Pour que WSL ``voie'' le
transcepteur, il faut attacher l'USB depuis PowerShell :

\begin{lstlisting}[style=shell, caption={Attacher l'USB du transcepteur a WSL (PowerShell admin)}]
usbipd list                       # voir le BUSID du CP210x (transcepteur)
usbipd bind --busid <BUSID>
usbipd attach --wsl --busid <BUSID>   # dans WSL apparait comme /dev/ttyUSBx
\end{lstlisting}

\textbf{Alternative --- Windows natif (sans WSL) :} installer Docker Desktop et
Python pour Windows et lancer les composants à la main : \texttt{docker compose up}
pour MongoDB, puis \texttt{python serial\_link.py} et \texttt{python webserver.py}.
Le port s\'erie sera \texttt{COMx} (à v\'erifier dans le Gestionnaire de
p\'eriph\'eriques) et il faut ajuster \texttt{SERIAL\_PORT\_REAL} dans le fichier
\file{.env} à ce \texttt{COMx}. Dans les deux cas, l'interface s'ouvre sur
\texttt{http://localhost:5000}.

\begin{warningbox}
  \textbf{R\'eseau WiFi --- EDUROAM ne fonctionne pas.} Comme l'IHM utilise MongoDB
  (dans Docker) et la pile r\'eseau locale, sur le r\'eseau institutionnel
  \textbf{EDUROAM l'IHM ne fonctionne pas} : EDUROAM applique un isolement des clients
  et des restrictions r\'eseau (pare-feu/proxy) qui emp\^echent le bon fonctionnement
  de Docker/MongoDB et de la communication locale n\'ecessaire.

  \textbf{Solution (test\'ee par le groupe) :} connecter l'ordinateur qui ex\'ecute
  l'IHM à un r\'eseau sans ces restrictions. En pratique, \textbf{partager les
  donn\'ees mobiles du t\'el\'ephone (hotspot / partage de connexion) vers
  l'ordinateur} : avec le hotspot du t\'el\'ephone, l'IHM d\'emarre et fonctionne
  correctement.
\end{warningbox}

% =============================================================================
\section{Pr\'esentation du projet}
% =============================================================================

\subsection{Concept physique}

Le drone est un v\'ehicule de surface propuls\'e principalement par une voile. Sa
conception est originale : la voile est à \textbf{profil d'aile} et tourne librement
autour d'un axe vertical solidaire de la coque. Son a\'erodynamique passive oriente
automatiquement le profil de sorte que le vent arrive toujours à environ
$45^{\circ}$ de l'axe longitudinal du drone, sans besoin d'un contr\^ole actif de
l'angle de la voile. La man\oe uvre se fait avec deux actionneurs principaux :

\begin{itemize}
  \item \textbf{Servo aileron (voile) --- Futaba S3003 :} contr\^ole un petit volet
    sur le bord de fuite de la voile. Sa course utile est strictement de
    $\pm10^{\circ}$. La d\'eflexion d\'etermine de quel c\^ot\'e (b\^abord ou tribord)
    la voile pousse le bateau, choisissant de fait l'amure sur laquelle il navigue.
    \textbf{Ce n'est pas un contr\^ole continu} : seules les deux positions extr\^emes
    ont un sens physique, donc on la mod\'elise comme un \'etat binaire ($+1$ / $-1$).

  \item \textbf{Servo safran / rotor --- Graupner Regatta ECO II (mod.\ 5176) :}
    c'est l'actionneur principal de cap (lacet). C'est un \emph{treuil de voile
    multi-tour positionnel} : le signal PWM commande une \textbf{position} de l'arbre
    (et non une vitesse). Centre m\'ecanique = \us{1500}, course totale $=$ 6 tours,
    et il maintient la position sous charge comme un servo standard.

  \item \textbf{Propulseur (ESC) :} utilis\'e uniquement en l'absence de vent
    suffisant. Permet la propulsion à h\'elice en secours. Sur le PCB V2 on utilise un
    seul ESC (le second a \'et\'e supprim\'e).
\end{itemize}

\begin{infobox}[Principe de navigation sans capteur de vent]
  Le syst\`eme \textbf{ne mesure pas} la direction du vent directement : il la
  \emph{d\'eduit} à partir du track GPS et de l'\'etat des actionneurs. L'estimation
  n\'ecessite un d\'eplacement minimal de 30~m. Ce manque de capteurs est central
  dans le projet et est repris en Section~\ref{sec:abiertos}.
\end{infobox}

\subsection{Objectif du projet}

Le projet est un PER (Projet d'\'Etude et de Recherche) dont le but final est
l'\textbf{acquisition autonome de donn\'ees marines} : le drone doit atteindre des
points de passage (\emph{waypoints}) GPS en naviguant à la voile, collecter des
donn\'ees de capteurs et les transmettre par radio LoRa à une station sol.

\subsection{Mat\'eriel principal}

\begin{longtable}{L{4cm} L{9cm}}
  \toprule
  \rowcolor{tableheadbg}
  \textcolor{tableheadfg}{\textbf{Composant}} &
  \textcolor{tableheadfg}{\textbf{Description}} \\
  \midrule
  \endhead
  \bottomrule
  \endfoot
  LilyGO T-Beam V1.1 & Module int\'egrant ESP32 WROVER, LoRa SX1276 433~MHz,
                       GPS u-blox NEO-6M, AXP192 (gestion d'\'energie) et
                       porte-batteries 18650 \\[3pt]
  Futaba S3003 & Servo aileron de la voile, centre à \us{1520}, $\pm10^{\circ}$ \\[3pt]
  Graupner Regatta ECO II & Treuil de safran multi-tour, \us{1500} = centre,
                            1000--2000~\textmu s = plage, jusqu'à 3,3~A en charge \\[3pt]
  Pro-Tronik PTR-6A & \'Emetteur 6 voies FHSS 2,4~GHz, PWM 1000--2000~\textmu s \\[3pt]
  Pro-Tronik R8X & R\'ecepteur 8 voies, PWM standard 50~Hz \\[3pt]
  Pro-Tronik Black Fet & ESC de la s\'erie, programmation par carte EPRG-3 \\[3pt]
  Antenne GPS & Taoglas active, connecteur u.FL \\[3pt]
  PCB v1 / PCB v2 & Carte d'interface sur mesure (KiCad 9) \\
\end{longtable}

% =============================================================================
\section{Code original --- \texttt{boat.ino}}
% =============================================================================

\subsection{Architecture et fonctionnement}

Le code original (\file{boat.ino}, du premier groupe, hors de ce d\'ep\^ot), \'ecrit
par Mohamed EL JILY, constitue le point de d\'epart du projet. Il \'etait organis\'e
en modules : \file{Gps}, \file{ServoControl}, \file{RadioReceiver},
\file{MotorControl} et la logique de navigation dans \file{boat.ino}. Il incluait
d\'ej\`a une logique de navigation à la voile (angle relatif vent/cap, changement
d'amure automatique, lofer/abattre) et la r\'eception de waypoints par LoRa en JSON.

\subsection{Probl\`emes et limitations identifi\'es}

À la reprise du projet, plusieurs probl\`emes bloquants ont \'et\'e identifi\'es qui,
conjointement avec le besoin de nouvelles fonctionnalit\'es discut\'ees avec le
client et le professeur, ont motiv\'e la r\'e\'ecriture compl\`ete du firmware.

\subsubsection{Biblioth\`eque ESP32Servo --- conflit LEDC / interruptions}

\begin{warningbox}
  La biblioth\`eque \texttt{ESP32Servo} utilise le p\'eriph\'erique \textbf{LEDC} de
  l'ESP32. LEDC est \textbf{incompatible} avec le d\'ecodage du PWM par interruptions
  GPIO utilis\'e pour lire la radio : lorsque LEDC est actif, les interruptions du
  r\'ecepteur RC sont perturb\'ees et les canaux lus \textbf{reviennent constamment à
  z\'ero}, m\^eme si les signaux d'entr\'ee sont corrects à l'oscilloscope. Cela rend
  impossible de contr\^oler les servos et de lire la radio en m\^eme temps. C'\'etait
  le probl\`eme le plus critique du projet.
\end{warningbox}

\subsubsection{Affectation des broches RC --- conflit I2C / AXP192}

\begin{warningbox}
  Le code original affectait des canaux RC à \gpio{21} et \gpio{22}, qui sont les
  broches I2C internes du T-Beam utilis\'ees par l'AXP192 (gestionnaire d'\'energie).
  De plus, \gpio{23} est reli\'e en interne au RESET du SX1276 (LoRa) : attacher une
  interruption sur cette broche pouvait provoquer des resets esp\'er\'es de la radio.
\end{warningbox}

\subsubsection{Broche batterie \gpio{35} --- drain de 20~mA continu}

Le T-Beam V1.1 poss\`ede un diviseur r\'esistif de \emph{faible imp\'edance}
($\approx370~\Omega$ total) soud\'e en permanence sur \gpio{35}. M\^eme sans lire
l'ADC, ce diviseur consomme $\approx20$~mA en continu ($\approx1{,}8$~W) directement
de la batterie. \textbf{Ne jamais lire \gpio{35}.}

\subsubsection{Broche LoRa RST --- erreur de documentation}

La documentation de LilyGO et de nombreux exemples indiquent \gpio{14} pour le RESET
du SX1276. Le hardware du T-Beam V1.1 c\^able physiquement le RESET sur \gpio{23}.

\subsubsection{Gestion de l'ESC --- sans armement ni limitation de variation}

Le code original n'impl\'ementait aucune fonction de l'ESC, car c'\'etait une
fonctionnalit\'e pr\'evue pour plus tard, tout comme la mesure de batterie.

% =============================================================================
\section{PCB Version 1 --- Carte d'interface originale}
% =============================================================================

La premi\`ere carte r\'ealisait l'interface entre le T-Beam et les actionneurs
(connecteurs 3 broches pour servos/ESC, connecteurs RC, r\'egulateur abaisseur 5~V
Pololu, bornes de puissance). Ses probl\`emes sont r\'esum\'es ci-dessous.

\begin{longtable}{L{3.5cm} L{9.5cm}}
  \toprule
  \rowcolor{tableheadbg}
  \textcolor{tableheadfg}{\textbf{Probl\`eme}} &
  \textcolor{tableheadfg}{\textbf{Impact}} \\
  \midrule
  \endhead
  \bottomrule
  \endfoot
  PWM 3,3~V sans amplification de niveau & Servos et ESC (logique 5~V) peuvent ne pas
    d\'etecter correctement les impulsions \\[3pt]
  Broches RC sur \gpio{21}/\gpio{22} & Conflit I2C/AXP192 \\[3pt]
  ESC sur \gpio{32}/\gpio{33} & Co\"incident avec LoRa DIO1/DIO2 du T-Beam \\[3pt]
  \gpio{35} pour la batterie & Diviseur de 370~$\Omega$ = 20~mA continus \\[3pt]
  Pas de protection de polarit\'e & Risque en cas d'alimentation incorrecte \\
\end{longtable}

% =============================================================================
\section{R\'e\'ecriture du firmware --- Architecture Post-Claude}
% =============================================================================

\subsection{Motivations}

La r\'e\'ecriture compl\`ete vise à corriger toutes les limitations identifi\'ees et
à pr\'eparer l'architecture pour les phases autonomes. Objectifs :

\begin{enumerate}
  \item Utiliser \textbf{MCPWM} au lieu de LEDC pour les sorties d'actionneurs.
  \item D\'ecodage RC \textbf{enti\`erement par interruptions} avec protection ISR.
  \item Architecture orient\'ee objet en couches, extensible par phases.
  \item Armement de l'ESC et limitation de variation (\emph{slew rate}).
  \item Gestion propre de l'\'energie via AXP192.
  \item GPS, LoRa et batterie fonctionnant simultan\'ement.
\end{enumerate}

\subsection{Architecture en couches}

\begin{lstlisting}[style=shell, caption={Arborescence du projet (racine = Post-Claude)}]
Post-Claude/                   <- dossier racine du projet
  main/main.ino                <- Point d'entree (begin / update)
  main/src/
    app/DroneApp.cpp           <- Orchestrateur (scheduler logiciel)
    config/BoardConfig.h       <- Affectation des GPIO et seuils RC
    config/Calibration.h       <- Valeurs en microsecondes (servo/ESC)
    core/Types.h               <- Types partages (RcFrame, ActuatorCommand)
    control/ModeManager        <- CH5 -> mode de controle
    control/ManualController   <- Mode manuel unifie (voile+safran+moteur)
    control/AutoController     <- Enveloppe de la navigation autonome
    drivers/RcReceiver         <- Lecture PWM par interruptions
    drivers/McpwmActuators     <- Sorties MCPWM (voile, safran, ESC)
    drivers/AxpPower           <- Init AXP192 (rails GPS/LoRa/3.3V)
    drivers/BatteryAdc         <- ADC tension batterie
    drivers/GpsUart            <- Lecture GPS (TinyGPSPlus + UBX)
    drivers/LoRaRadio          <- Driver SX1276
    comm/LoRaComm              <- Protocole JSON LoRa (heartbeat + commandes)
    navigation/                <- navigation.h, Navigator, MissionManager...
  transceiver/transceiver.ino  <- Station sol (pont USB <-> LoRa)
  IHM/                         <- Interface web de station sol (FastAPI)
\end{lstlisting}

La r\'epartition suit une logique claire : \emph{drivers} (mat\'eriel),
\emph{services / comm} (donn\'ees), \emph{navigation} (g\'eom\'etrie et mission),
\emph{control} (d\'ecisions) et \emph{app} (orchestration). L'ordonnanceur est
coop\'eratif (sans RTOS) : \file{DroneApp::update()} d\'eclenche chaque t\^ache
quand sa p\'eriode est \'echue, \'evitant \texttt{delay()} qui bloquerait la lecture
RF.

\subsection{D\'ecisions techniques fondamentales}

\subsubsection{MCPWM obligatoire --- jamais LEDC}

\begin{lstlisting}[style=cpp, caption={McpwmActuators.cpp -- initialisation MCPWM}]
bool McpwmActuators::begin() {
    mcpwm_gpio_init(MCPWM_UNIT_0, MCPWM0A, BoardConfig::SAIL_SERVO_PIN);
    mcpwm_gpio_init(MCPWM_UNIT_0, MCPWM0B, BoardConfig::ROTOR_SERVO_PIN);
    mcpwm_gpio_init(MCPWM_UNIT_0, MCPWM1A, BoardConfig::ESC1_PIN);
    mcpwm_config_t cfg = {};
    cfg.frequency    = Calibration::PWM_FREQ_HZ;   // 50 Hz
    cfg.duty_mode    = MCPWM_DUTY_MODE_0;          // actif haut
    cfg.counter_mode = MCPWM_UP_COUNTER;
    mcpwm_init(MCPWM_UNIT_0, MCPWM_TIMER_0, &cfg);  // voile + safran
    mcpwm_init(MCPWM_UNIT_0, MCPWM_TIMER_1, &cfg);  // ESC
    ...
}
\end{lstlisting}

Le p\'eriph\'erique MCPWM est ind\'ependant des interruptions GPIO : les deux
coexistent sans se perturber. Cette contrainte est \textbf{non n\'egociable} dans ce
projet.

\subsubsection{Interruptions RC avec protection portMUX}

\begin{lstlisting}[style=cpp, caption={RcReceiver.cpp -- ISR et section critique}]
void IRAM_ATTR RcReceiver::handleEdge(uint8_t pin, Ch ch) {
    const uint32_t now = micros();
    const int      idx = static_cast<int>(ch);
    if (gpio_get_level(static_cast<gpio_num_t>(pin))) {
        portENTER_CRITICAL_ISR(&mux_);
        riseUs_[idx] = now;                        // front montant
        portEXIT_CRITICAL_ISR(&mux_);
    } else {
        uint32_t rise;
        portENTER_CRITICAL_ISR(&mux_);
        rise = riseUs_[idx];
        portEXIT_CRITICAL_ISR(&mux_);
        const uint32_t width = now - rise;         // largeur d'impulsion
        if (width >= RC_VALID_MIN_US && width <= RC_VALID_MAX_US) {
            portENTER_CRITICAL_ISR(&mux_);
            pulseUs_[idx]     = static_cast<uint16_t>(width);
            lastValidUs_[idx] = now;
            portEXIT_CRITICAL_ISR(&mux_);
        }
    }
}
\end{lstlisting}

Les ISR sont plac\'ees en RAM (\texttt{IRAM\_ATTR}) et se synchronisent avec la
boucle principale via \texttt{portMUX} (m\'ecanisme FreeRTOS appropri\'e pour l'ESP32
dual-core). Un \emph{timeout} de 100~ms marque un canal comme perdu.

\subsubsection{Toutes les consignes d'actionneurs en microsecondes}

Tous les actionneurs sont command\'es exclusivement en microsecondes, jamais en
degr\'es. Cela \'elimine une couche de conversion source d'erreurs et correspond
directement au protocole PWM des servos et ESC.

\subsubsection{Slew rate --- protection contre les à-coups}

\begin{lstlisting}[style=cpp, caption={McpwmActuators.cpp -- limitation de variation}]
uint16_t McpwmActuators::slew(uint16_t current, uint16_t target, uint16_t step) {
    if (current < target) {
        uint32_t next = (uint32_t)current + step;
        return (next >= target) ? target : (uint16_t)next;
    }
    if (current > target)
        return (current > target + step) ? (uint16_t)(current - step) : target;
    return current;
}
// ESC_SLEW_US = 30 us/tick, appele toutes les 20 ms -> max 1500 us/s
\end{lstlisting}

% =============================================================================
\section{Phase 1 --- Contr\^ole manuel par radio}
% =============================================================================

\subsection{R\'ealisation}

La Phase~1 couvre le contr\^ole manuel complet. On a impl\'ement\'e les drivers
(\file{AxpPower}, \file{RcReceiver}, \file{McpwmActuators}), le \file{ModeManager},
le \file{ManualController} et l'orchestrateur \file{DroneApp} (ticks de 20~ms pour
le contr\^ole, 200~ms pour le debug, 1~s pour LoRa).

\subsection{Modes de vol (s\'electeur CH5) --- version unifi\'ee actuelle}

Le s\'electeur 3 positions est c\^abl\'e sur CH5. La r\'epartition des modes a \'et\'e
\textbf{r\'eorganis\'ee} par rapport à la version initiale (qui avait deux modes
manuels s\'epar\'es ``ManualServo'' et ``ManualProp'') :

\begin{longtable}{C{2.6cm} C{2.4cm} L{8cm}}
  \toprule
  \rowcolor{tableheadbg}
  \textcolor{tableheadfg}{\textbf{CH5}} &
  \textcolor{tableheadfg}{\textbf{Mode}} &
  \textcolor{tableheadfg}{\textbf{Comportement}} \\
  \midrule
  \endhead
  \bottomrule
  \endfoot
  $\leq$ \us{1250} & Automatic & Navigation autonome (suit la mission LoRa) \\[3pt]
  1400--1600 & Sail & \textbf{Inerte} : toutes les sorties au neutre \\[3pt]
  $>$ \us{1800} & Manual & Mode manuel \textbf{unifi\'e} (voir ci-dessous) \\[3pt]
  Signal perdu & Failsafe & Suit la mission autonome (comme Automatic) \\
\end{longtable}

\begin{warningbox}
  \textbf{Attention s\'ecurit\'e :} en cas de perte RC, le syst\`eme passe en
  \emph{Failsafe}, qui se comporte comme le mode autonome (navigue seul). Comme le
  mode autonome n'a jamais \'et\'e test\'e sur l'eau, le \textbf{premier essai r\'eel
  doit \^etre manuel} et avec un plan de r\'ecup\'eration.
\end{warningbox}

\subsection{Mode manuel unifi\'e}

Un seul mode (\texttt{computeManual}) contr\^ole à la fois la voile, le safran et le
propulseur :

\begin{itemize}
  \item \textbf{CH2 $\to$ voile (binaire) :} bande morte $\pm35$~\textmu s autour de
    \us{1500}. Le premier frame initialise l'\'etat depuis la position du manche pour
    \'eviter un saut à l'entr\'ee du mode. Dans la bande morte, le dernier \'etat est
    maintenu.
  \item \textbf{CH4 $\to$ safran (positionnel) :} la course mesur\'ee de CH4
    (1180--1790~\textmu s) est mapp\'ee sur la plage physique du treuil
    (\texttt{ROTOR\_MIN/MAX\_US} = 1417--1583~\textmu s, $\pm90^{\circ}$). Une bande
    morte centrale fixe le centre exact (\us{1500}).
  \item \textbf{CH3 $\to$ propulseur (bidirectionnel, voir Section~\ref{sec:bidir}).}
\end{itemize}

\begin{lstlisting}[style=cpp, caption={ManualController -- logique de voile binaire avec etat memorise}]
if (!sailStateInitialized_) {
    sailState_ = (frame.ch2 >= RC_MID_US) ? +1 : -1;
    sailStateInitialized_ = true;
}
if (frame.ch2 > RC_MID_US + db)      sailState_ = +1;
else if (frame.ch2 < RC_MID_US - db) sailState_ = -1;
cmd.sailUs = (sailState_ > 0) ? SAIL_PLUS_US : SAIL_MINUS_US;
\end{lstlisting}

\subsection{Propulseur bidirectionnel (marche avant / arri\`ere)}
\label{sec:bidir}

Dans la version initiale, le moteur n'allait que de 0 à 100~\% (commande invers\'ee,
avec l'ESC en mode unidirectionnel, \us{1000} = arr\^et). On l'a modifi\'e pour que
le moteur aille de \textbf{$+100~\%$ (avant) à $-100~\%$ (arri\`ere)} en utilisant
l'ESC en mode \textbf{bidirectionnel}, o\`u le neutre passe à \us{1500} :

\begin{longtable}{C{4.5cm} C{3cm} L{5cm}}
  \toprule
  \rowcolor{tableheadbg}
  \textcolor{tableheadfg}{\textbf{Position du manche CH3}} &
  \textcolor{tableheadfg}{\textbf{Puissance}} &
  \textcolor{tableheadfg}{\textbf{ESC}} \\
  \midrule
  \endhead
  \bottomrule
  \endfoot
  \texttt{CH3\_FULL\_US} (\us{1100}) & $+100~\%$ avant & \us{2000} \\[3pt]
  \texttt{CH3\_CENTER\_US} (\us{1545}) $\pm$ bande morte & 0~\% arr\^et & \us{1500} \\[3pt]
  \texttt{CH3\_ZERO\_US} (\us{1990}) & $-100~\%$ arri\`ere & \us{1000} \\
\end{longtable}

Comme le manche des gaz est de type \emph{cran} (il ne se recentre pas),
l'arr\^et est plac\'e au centre de sa course avec une bande morte
(\texttt{ESC\_CENTER\_DEADBAND\_US} $=40$~\textmu s). Pour garder tout le firmware
coh\'erent, le neutre \us{1500} est aussi utilis\'e au d\'emarrage, en Failsafe, dans
le mode Sail et comme arr\^et du mode autonome~; la limite inf\'erieure de saturation
a \'et\'e abaiss\'ee à \us{1000} pour permettre la marche arri\`ere.

\begin{warningbox}
  La marche arri\`ere ne fonctionne physiquement que si l'ESC est \textbf{programm\'e
  en mode bidirectionnel} avec la carte EPRG-3. Sinon, toute valeur en dessous de
  \us{1500} est ignor\'ee et on n'obtient que la marche avant.
\end{warningbox}

\subsection{R\'esultats des tests --- Phase 1}

\begin{itemize}
  \item \done\ L'aileron bascule correctement entre \us{1465} et \us{1575}.
  \item \done\ Le safran r\'epond proportionnellement au manche CH4.
  \item \done\ Le Failsafe se d\'eclenche en moins de 100~ms apr\`es la perte du signal.
  \item \done\ Compilation \texttt{arduino-cli} : $\approx$28~\% flash, 7~\% RAM.
  \item \todo\ Calibration de banc de \texttt{SAIL\_PLUS/MINUS\_US} et de
    \texttt{CH3\_CENTER\_US} en attente.
\end{itemize}

% =============================================================================
\section{Phase 2 --- GPS et navigation}
% =============================================================================

On a impl\'ement\'e \file{GpsUart} (Serial1 sur \gpio{34}/\gpio{12}, TinyGPSPlus),
\file{Navigator} (haversine : distance et cap), \file{MissionPlan} (jusqu'à 16
waypoints) et \file{MissionManager} (machine à \'etats
Idle$\to$Running$\to$Returning$\to$Complete). Deux probl\`emes critiques de GPS ont
\'et\'e r\'esolus dans cette phase.

\subsection{Probl\`eme critique --- GPS muet apr\`es configuration en flash}

\begin{warningbox}
  \textbf{Sympt\^ome :} le compteur \texttt{chars} reste bloqu\'e, le GPS est
  aliment\'e mais ne produit aucune phrase NMEA.

  \textbf{Cause :} une session pr\'ec\'edente (IDE Arduino ou u-center) a sauvegard\'e
  dans la flash interne du NEO-6M une configuration CFG-PRT qui d\'esactive la sortie
  NMEA. Cette configuration survit aux coupures d'alimentation.
\end{warningbox}

\begin{okbox}
  Envoyer un reset usine \texttt{UBX CFG-CFG} au d\'emarrage (efface la configuration
  flash et charge les valeurs ROM). Impl\'ement\'e dans \file{GpsUart::begin()}.
\end{okbox}

\subsection{Probl\`eme critique --- antenne active non aliment\'ee}
\label{subsec:antena}

\begin{warningbox}
  \textbf{Sympt\^ome :} NMEA re\c{c}u mais \texttt{visible = 0} satellites m\^eme en
  ext\'erieur d\'egag\'e et avec antenne externe connect\'ee.

  \textbf{Cause :} le superviseur d'antenne du NEO-6M est d\'esactiv\'e par d\'efaut.
  Sans l'activer, aucune tension de polarisation n'est fournie au LNA de l'antenne
  active et le commutateur RF du T-Beam reste sur l'antenne c\'eramique interne,
  court-circuitant le connecteur u.FL.
\end{warningbox}

\begin{okbox}
  Envoyer \texttt{UBX CFG-ANT} avec \texttt{flags = 0x001B} (contr\^ole de tension RF,
  d\'etection de court-circuit, r\'ecup\'eration). Apr\`es la correction : fix en
  moins de 60~s, 7 satellites, HDOP $\approx$2,2.
\end{okbox}

\subsection{Note sur le d\'emarrage à froid}

Sans batterie LiPo connect\'ee, chaque allumage perd l'almanach GPS : le d\'emarrage
à froid prend $\approx$2~min en ext\'erieur. Avec la batterie, les \'eph\'em\'erides
sont conserv\'ees et le fix s'obtient en quelques secondes (d\'emarrage à chaud).

% =============================================================================
\section{Phase 3 --- LoRa et t\'el\'em\'etrie}
% =============================================================================

On a impl\'ement\'e le driver \file{LoRaRadio} (SX1276, SPI HSPI, 433~MHz), la couche
de protocole \file{LoRaComm} (heartbeat JSON à 1~Hz + dispatcher de commandes) et le
sketch de station sol \file{transceiver.ino} (pont USB$\leftrightarrow$LoRa avec CLI
compacte). Tout a \'et\'e v\'erifi\'e sur mat\'eriel (RSSI $-101$ à $-106$~dBm à
courte distance).

\subsection{Protocole JSON --- heartbeat actuel}

Le format a \'et\'e affin\'e pour rester largement sous la limite de 255~octets de la
FIFO du SX1276 (champs redondants supprim\'es, indicateurs de sant\'e GPS/RC
ajout\'es, et le RSSI est inject\'e par le transcepteur au sol) :

\begin{lstlisting}[style=json, caption={Heartbeat drone -> sol (1 Hz)}]
{"origin":"boat","type":"info","message":{
  "mode":"standby|route-ready|navigate",
  "location":[48.36000,-4.56000],
  "servos":{"sail":10,"rudder":-12},
  "heading":270,"wind":180,"bat":7.42,
  "fix":1,"sat":7,"hdop":2.2,"rc":1,
  "wt":3,"wc":1,"wobs":42,"rssi":-104}}
\end{lstlisting}

\begin{itemize}
  \item \texttt{servos.sail} $=\pm10^{\circ}$ (binaire)~; \texttt{servos.rudder} en
    degr\'es du treuil. \texttt{fix/sat/hdop} = sant\'e GPS~; \texttt{rc}=1 si le
    r\'ecepteur d\'elivre des impulsions.
  \item \texttt{wt/wc} = waypoints total/actuel. \texttt{wobs} = progression de la
    mesure de vent (0--100~\%), pr\'esent \emph{seulement} pendant la mesure.
  \item \texttt{rssi} n'est \textbf{pas} envoy\'e par le bateau : il est inject\'e par
    le transcepteur au sol, il ne compte donc pas dans le budget de 255~octets.
\end{itemize}

\subsection{Commandes (sol $\to$ drone)}

\begin{lstlisting}[style=json, caption={Commandes support\'ees}]
{"origin":"server","type":"command","message":{"navigate":true}}
{"origin":"server","type":"command","message":{"stop":true}}
{"origin":"server","type":"command","message":{"restart":true}}
{"origin":"server","type":"command","message":{"wind-observation":true}}
{"origin":"server","type":"command","message":{"wind-command":{"value":225}}}
{"origin":"server","type":"command","message":{"home":{"lat":48.36,"lon":-4.56}}}
{"origin":"server","type":"command","message":
  {"waypoints":{"number":2,"points":"48.36,-4.56,48.37,-4.55"}}}
\end{lstlisting}

\begin{infobox}[Note sur le TX bloquant]
  \texttt{LoRa.endPacket()} est synchrone. À SF7/BW125~kHz, un paquet de
  $\approx$210~octets prend $\approx$450~ms, bloquant la boucle principale une fois
  par seconde. Acceptable pour l'instant~; il faudra passer en TX asynchrone avant la
  navigation autonome sur l'eau si le jitter de la boucle de contr\^ole devient
  critique.
\end{infobox}

% =============================================================================
\section{PCB Version 2 --- Am\'eliorations mat\'erielles}
% =============================================================================

La V2 du PCB corrige tous les probl\`emes de la V1 et int\`egre les enseignements du
d\'eveloppement logiciel.

\subsection{Correspondance des broches (PCB V2)}

\begin{longtable}{C{3cm} C{1.6cm} L{7.5cm}}
  \toprule
  \rowcolor{tableheadbg}
  \textcolor{tableheadfg}{\textbf{Signal}} &
  \textcolor{tableheadfg}{\textbf{GPIO}} &
  \textcolor{tableheadfg}{\textbf{Notes}} \\
  \midrule
  \endhead
  \bottomrule
  \endfoot
  RC CH2 (aileron) & 39 & SVN, entr\'ee seule \\[2pt]
  RC CH3 (gaz) & 14 & Lib\'er\'e (ne partage plus l'I2C) \\[2pt]
  RC CH4 (safran) & 13 & Lib\'er\'e \\[2pt]
  RC CH5 (mode) & 4 & Commutateur 3 positions \\[2pt]
  Servo voile (PWM1) & 2 & Via transistor BC548 \\[2pt]
  Servo safran (PWM2) & 25 & Via transistor BC548 \\[2pt]
  ESC1 & 15 & Via transistor BC548 (ESC unique) \\[2pt]
  ADC batterie & 36 & SVP, ADC1\_CH0, entr\'ee seule \\[2pt]
  LoRa SCK/MISO/MOSI/CS & 5/19/27/18 & SPI HSPI \\[2pt]
  LoRa RST & 23 & Maintenant libre (CH6 supprim\'e) \\[2pt]
  LoRa IRQ (DIO0) & 26 & \\[2pt]
  GPS RX / TX & 34 / 12 & Serial1 \\[2pt]
  I2C SDA / SCL & 21 / 22 & AXP192 interne + \textbf{connecteur d'extension} \\
\end{longtable}

\subsection{Amplification de niveau logique --- transistors BC548}
\label{subsec:bc548}

Les sorties PWM de l'ESP32 sont à 3,3~V~; les servos attendent 5~V. La V2 utilise des
transistors NPN BC548 en \'emetteur commun.

\begin{infobox}[D\'ecouverte importante --- DUTY\_MODE\_0 sur les deux canaux]
  Initialement, le canal du safran (OPR\_B) utilisait \texttt{MCPWM\_DUTY\_MODE\_1}
  (inversion), en supposant que le NPN inversait le signal et qu'il fallait
  pr\'e-inverser. En r\'ealit\'e, le montage BC548 inverse d\'ej\`a~; appliquer
  \texttt{DUTY\_MODE\_1} produisait une \textbf{double inversion} qui rendait le servo
  du safran inop\'erant. La solution est d'utiliser \texttt{DUTY\_MODE\_0} (actif
  haut) sur \textbf{les deux} canaux.
\end{infobox}

\subsection{Mesure de batterie --- diviseur haute imp\'edance}

On utilise un diviseur R5 $=562$~k$\Omega$ / R6 $=120$~k$\Omega$ (rapport $5{,}683$)
sur \gpio{36} avec une att\'enuation de 11~dB. À 7,4~V le diviseur consomme
$\approx10{,}9$~\textmu A, n\'egligeable.

\begin{warningbox}
  \textbf{D\'etail de montage :} sur une carte, les r\'esistances ont \'et\'e
  soud\'ees invers\'ees (120~k en haut, 562~k en bas), donnant une lecture de 17,9~V
  au lieu de $\approx$5~V. Le code \'etait correct~; la solution a \'et\'e d'inverser
  physiquement les r\'esistances.
\end{warningbox}

\subsection{AXP192 --- I2C persistant}

Dans la V1, on appelait \texttt{Wire.end()} apr\`es l'init de l'AXP192 pour lib\'erer
\gpio{21}/\gpio{22}. Dans la V2, ces broches restent libres pour l'I2C, donc
\texttt{Wire.end()} est supprim\'e : le bus I2C reste actif et disponible pour de
futurs capteurs sur le connecteur d'extension.

% =============================================================================
\section{Navigation autonome (Phases 5 et 6)}
% =============================================================================

\subsection{Phase 5 --- Estimation du vent à partir du track GPS}

Sans capteur de vent, la direction est infer\'ee en naviguant sur une amure fixe : le
bateau fixe la voile à $+10^{\circ}$, lisse le cap GPS avec une moyenne mobile
circulaire et, apr\`es avoir parcouru \texttt{WIND\_OBS\_DISTANCE\_M} $=30$~m, fixe
le vent à \texttt{cap\_liss\'e $+ 90^{\circ}$}. Il existe aussi une commande LoRa
\texttt{wind-command} pour forcer la direction manuellement (secours fiable).

\begin{warningbox}
  Le d\'ecalage de $+90^{\circ}$ et le seuil de 30~m n\'ecessitent une validation sur
  le terrain. De plus, pendant la mesure, le heartbeat est identique au repos sauf
  pour le champ \texttt{wobs}.
\end{warningbox}

\subsection{Phase 6 --- Algorithme de navigation à voile}

L'algorithme de navigation (\file{navigation.h}) a \'et\'e d\'evelopp\'e par un membre
du groupe (branche \texttt{Testautoboat} du d\'ep\^ot
\texttt{DeltaGod/PER\_MEA\_LYINCROYAP}) et int\'egr\'e à ce firmware. Il impl\'emente :

\begin{itemize}
  \item Zones interdites au pr\`es et au portant (\emph{upwind/downwind}) avec zigzag
    de louvoyage.
  \item \'Evitement d'empannage (\emph{empannage}) via une machine à \'etats.
  \item Couloir de tol\'erance lat\'erale (\emph{cross-track}) autour de la ligne vers
    le waypoint.
  \item Lofer/abattre (\emph{lofer/abattre}) pour le suivi fin du cap.
\end{itemize}

La derni\`ere version int\'egr\'ee apporte : une s\'election d'amure initiale plus
intelligente (\texttt{nav\_sideMovingTowardWaypoint}, qui choisit le c\^ot\'e qui
progresse le plus vers le waypoint), une d\'etection d'arriv\'ee par ``plan du
waypoint d\'epass\'e à l'int\'erieur du couloir'' (\texttt{nav\_hasPassedWaypoint}) et
un couloir par d\'efaut \'elargi de 20 à 100~m. L'enveloppe \file{AutoController}
mappe le cap de safran de l'algorithme 1:1 sur les degr\'es physiques du treuil,
limit\'e à $\pm20^{\circ}$ (\texttt{ROTOR\_AUTO\_RANGE\_DEG}), car le $\pm90^{\circ}$
complet se r\'ev\'elait trop agressif.

\begin{warningbox}
  La navigation autonome est \textbf{cod\'ee mais jamais test\'ee sur l'eau} et n'a
  pas de tests unitaires. Il manque aussi un \emph{WinchTracker} : le Regatta ECO II
  n'a pas de retour de position, donc la position r\'eelle du safran n'est pas connue
  (voir Section~\ref{sec:abiertos}).
\end{warningbox}

% =============================================================================
\section{Interface de station sol (IHM)}
% =============================================================================

La station sol se compose d'un second T-Beam avec \file{transceiver.ino}
(pont USB$\leftrightarrow$LoRa) et d'une application web (\file{IHM/}) en Python
FastAPI + MongoDB + une carte Leaflet. Fonctions principales :

\begin{itemize}
  \item Carte avec la position du bateau, waypoints (clic pour les d\'efinir), envoi
    de route et \emph{polling} de t\'el\'em\'etrie à 1~Hz.
  \item Panneau de sant\'e avec voyants (Liaison, GPS, Batterie, RC, Contr\^ole) et RSSI.
  \item Boussole de vent et pr\'ediction de trajectoire (r\'eplique en JavaScript de
    \file{navigation.h}).
  \item Barre de progression de la mesure de vent.
  \item Autod\'etection du port du transcepteur par num\'ero de s\'erie du CP2104
    (transcepteur $=$ \texttt{01C00B54}, bateau $=$ \texttt{02126CF1}).
\end{itemize}

% =============================================================================
\section{Probl\`emes r\'esolus et solutions d\'etaill\'ees}
% =============================================================================

Cette section r\'eunit, de mani\`ere exhaustive, les probl\`emes significatifs
rencontr\'es pendant le d\'eveloppement et comment ils ont \'et\'e r\'esolus.

\subsection{Alimentation du T-Beam --- impossible d'alimenter par les broches}
\label{sec:alim}

\begin{warningbox}
  \textbf{Probl\`eme :} le T-Beam V1.1 \textbf{ne peut pas \^etre aliment\'e en
  injectant 5~V par les broches du header d'extension}. Les seules voies
  d'alimentation valides sont le \textbf{porte-batteries 18650/LiPo} (connecteur JST
  de la carte elle-m\^eme) et le \textbf{port USB}. Le PCB V2, en revanche, a \'et\'e
  con\c{c}u pour alimenter le T-Beam via la broche \texttt{+5V} du header d'extension
  (sortie du r\'egulateur abaisseur Pololu).

  La raison est la topologie interne d'\'energie du T-Beam : la broche \texttt{+5V}
  du header se trouve sur le chemin du VBUS de l'USB, apr\`es la diode de protection,
  et \textbf{n'entre pas de fa\c{c}on fiable dans l'AXP192} (le PMIC qui active les
  rails de 3,3~V, GPS et LoRa). Par cons\'equent, aliment\'e seulement par cette
  broche, l'AXP192 \textbf{ne d\'emarre} souvent \textbf{pas}, et sous charge la
  tension s'effondre (on a observ\'e le bateau tournant à $\approx$0,81~V, clairement
  en \emph{brownout}), ce qui provoque en plus la r\'e-\'enum\'eration continue du
  convertisseur USB-s\'erie (messages noyau \texttt{cp210x ... status -19}).
\end{warningbox}

\begin{okbox}
  \textbf{Solution rapide adopt\'ee :} on a coup\'e un c\^able USB-A~$\to$~USB-C~; le
  bout USB-A (lignes \texttt{+5V} et \texttt{GND}) a \'et\'e \textbf{soud\'e
  directement au rail de 5~V du PCB} (sortie du Pololu), et le bout USB-C se connecte
  au port USB du T-Beam. Ainsi, l'\'energie du PCB entre dans le T-Beam par son
  \textbf{chemin d'alimentation USB valide} (VBUS $\to$ AXP192), au lieu de la broche
  \texttt{+5V} du header, qui n'arrivait pas à d\'emarrer le PMIC.

  \textbf{Recommandation :} cette solution est un \emph{patch} fonctionnel mais peu
  \'el\'egant (un c\^able USB soud\'e au rail). Une correction d\'efinitive dans une
  future r\'evision du PCB serait de router l'alimentation du PCB directement vers
  l'entr\'ee correcte du T-Beam (VBUS ou le JST de batterie), en \'evitant la broche
  \texttt{+5V} du header.
\end{okbox}

\subsection{Moteur qui fonctionne ``par impulsions'' --- coupure par basse tension}
\label{subsec:lvc}

\begin{warningbox}
  \textbf{Sympt\^ome :} avec la batterie à $\approx$7~V (selon la t\'el\'em\'etrie), le
  moteur ne tourne pas de fa\c{c}on continue : il d\'emarre un instant, se coupe, et
  ne repart que si l'on ram\`ene la commande à 0~\% puis de nouveau à 100~\% un peu
  plus tard.

  \textbf{Cause (ce n'est pas un bug logiciel) :} c'est la \textbf{coupure par basse
  tension (LVC)} de l'ESC. Une batterie 2S à 7~V est d\'ej\`a basse
  ($\approx$3,5~V/cellule)~; sous la charge du moteur, la tension passe sous le seuil
  de coupure de l'ESC ($\approx$3,0--3,2~V/cellule) et l'ESC coupe le moteur pour
  prot\'eger le pack. Sans charge, la tension se r\'etablit, la coupure se r\'earme et
  le moteur retourne tourner un instant, produisant justement ce comportement ``par
  impulsions''. L'effet empire beaucoup si l'ESC a un mauvais r\'eglage du nombre de
  cellules.
\end{warningbox}

\begin{okbox}
  \textbf{Solution :} cha\^iner en s\'erie 2 batteries LiPo ($\approx$15,4~V). Aucun
  changement de code ne corrige une coupure par basse tension~; le firmware envoie une
  consigne stable.
\end{okbox}

\subsection{\'Echec de programmation du bateau --- brownout pendant le flash}

\begin{warningbox}
  \textbf{Sympt\^ome :} la programmation \'echoue r\'ep\'etitivement (``\emph{chip
  stopped responding}'', ``\emph{No serial data received}'') et le CP2104 se
  r\'e-\'enum\`ere au milieu de l'\'ecriture. Cela n'arrive que lorsque le T-Beam est
  connect\'e directement au PCB.

  \textbf{Cause :} le m\^eme \emph{brownout} d'alimentation. Le pic de courant de
  l'effacement/\'ecriture de la flash effondre la tension et l'ESP32 se r\'einitialise.
\end{warningbox}

\begin{okbox}
  D\'econnecter le T-Beam du PCB, le programmer et, sans le d\'econnecter de l'USB, le
  replacer sur le PCB.
\end{okbox}

\subsection{Autres probl\`emes de firmware d\'ej\`a r\'esolus}

\subsubsection{DUTY\_MODE\_1 sur OPR\_B --- servo du safran inop\'erant}

Double inversion à travers le BC548 (voir Section~\ref{subsec:bc548}). Solution :
\texttt{DUTY\_MODE\_0} sur les deux canaux.

\subsubsection{Seuil d'armement de l'ESC trop bas}

Le minimum r\'eel du manche des gaz du PTR-6A est $\approx$\us{1265}, donc le seuil
d'armement de \us{1050} n'\'etait jamais atteint. Il a \'et\'e port\'e à \us{1300}.
\textbf{Note :} le geste d'armement a \'et\'e \textbf{supprim\'e} ensuite, au passage
au mode manuel unifi\'e et au contr\^ole bidirectionnel (il \'etait incompatible avec
le nouveau mapping des gaz).

\subsubsection{Mode bloqu\'e en Automatic --- confusion CH4/CH5}

La version initiale lisait CH4 pour le mode, alors que le commutateur est sur CH5.
Solution : lire \texttt{frame.ch5}.

\subsubsection{Port s\'erie occup\'e lors de la programmation}

Un moniteur s\'erie ou un script Python gardait le port ouvert
(\texttt{Errno 16: Device or resource busy}). Solution : fermer le moniteur ou
\texttt{fuser /dev/ttyUSB0} et tuer le processus.

% =============================================================================
\section{Probl\`emes rencontr\'es --- options d'am\'elioration}
\label{sec:abiertos}
% =============================================================================

Cette section r\'eunit les probl\`emes encore \textbf{ouverts} et, pour chacun,
l'option d'am\'elioration propos\'ee pour le r\'esoudre. Ce sont les axes de travail
prioritaires pour que le drone puisse r\'eellement accomplir sa mission.

\subsection{Port\'ee LoRa tr\`es faible à l'int\'erieur du bo\^itier \'electronique}

\begin{warningbox}
  \textbf{Probl\`eme :} avec l'antenne LoRa actuelle plac\'ee \textbf{à l'int\'erieur
  du bo\^itier \'electronique}, la port\'ee radio est tr\`es faible. Le bo\^itier fait
  \'ecran et la proximit\'e de l'\'electronique d\'egrade le diagramme de rayonnement,
  r\'eduisant drastiquement le RSSI et la port\'ee utile de t\'el\'em\'etrie et de
  commandes.
\end{warningbox}

\begin{okbox}
  \textbf{Am\'elioration propos\'ee :} acqu\'erir une \textbf{antenne LoRa externe}
  (433~MHz) et la placer \textbf{sur le m\^at principal}, le plus haut et d\'egag\'e
  possible, hors du bo\^itier. Cela am\'eliore la ligne de vue et \'eloigne l'antenne
  des masses m\'etalliques et du bruit de l'\'electronique.

  \textbf{À consid\'erer --- connecteur :} il faut v\'erifier le connecteur de
  l'antenne et du module. Le SX1276 du T-Beam expose g\'en\'eralement
  \textbf{u.FL/IPEX}, tandis que les antennes externes 433~MHz ont habituellement du
  \textbf{SMA}. Il faudra le \textbf{c\^able adaptateur / pigtail correspondant
  (u.FL~$\to$~SMA)} et veiller à ce qu'il soit en 433~MHz (pas en 2,4~GHz) pour ne pas
  d\'egrader l'adaptation.
\end{okbox}

\subsection{Incoh\'erences du GPS}

\begin{warningbox}
  \textbf{Probl\`eme :} le GPS pr\'esente des incoh\'erences (fix intermittent, peu de
  satellites, position peu stable), aggrav\'ees par l'emplacement de l'antenne dans ou
  pr\`es du bo\^itier et de l'\'electronique.
\end{warningbox}

\begin{okbox}
  \textbf{Am\'elioration propos\'ee :} acqu\'erir \textbf{une autre antenne externe,
  cette fois de GPS}, et la placer \textbf{sur le pont}, avec une vue d\'egag\'ee du
  ciel et loin de l'antenne LoRa et de l'\'electronique.

  \textbf{À consid\'erer --- connecteur :} comme pour LoRa, il faut v\'erifier le
  connecteur. L'entr\'ee d'antenne GPS du NEO-6M est en \textbf{u.FL}~; les antennes
  GPS actives viennent souvent en \textbf{SMA} ou directement avec un c\^able u.FL. Il
  faudra pr\'evoir l'\textbf{adaptateur ad\'equat} et confirmer que l'antenne est
  \textbf{active} (le superviseur d'antenne est d\'ej\`a activ\'e par logiciel, voir
  Section~\ref{subsec:antena}).
\end{okbox}

\subsection{Sans mesure du vent ni de la position de la voile}

\begin{warningbox}
  \textbf{Probl\`eme :} le sch\'ema actuel \textbf{ne mesure pas d'o\`u vient le vent}
  ni \textbf{la position r\'eelle de la voile}. Comme le treuil du safran (Regatta ECO
  II) n'a pas non plus de retour de position, \textbf{il est impossible de conna\^itre
  la position r\'eelle du safran} ni de fermer la boucle de contr\^ole avec des
  donn\'ees r\'eelles : toute la navigation autonome d\'epend d'estimations (vent
  d\'eduit du track GPS, position de safran suppos\'ee à partir de la consigne).
\end{warningbox}

\begin{okbox}
  \textbf{Am\'elioration propos\'ee :} se procurer \textbf{2 capteurs} qui mesurent
  respectivement la \textbf{direction du vent} et la \textbf{position de la voile}, et
  les connecter au \textbf{bus I2C de la carte} via le \textbf{connecteur d\'ej\`a
  pr\'evu sur le PCB} (\gpio{21}/\gpio{22}, que le firmware a laiss\'es libres en
  supprimant \texttt{Wire.end()}). Avec la direction r\'eelle du vent et la position
  r\'eelle de la voile, il serait alors possible de d\'eduire la position effective du
  safran et de fermer correctement la boucle de navigation.

  \textbf{Avantage architectural :} le firmware r\'eserve d\'ej\`a le bus I2C et le PCB
  expose d\'ej\`a le connecteur d'extension, donc l'int\'egration \'electrique est
  directe~; il manquerait le driver \file{SensorBus} c\^ot\'e logiciel.
\end{okbox}

\subsection{Peu d'espace pour les batteries dans le bo\^itier}

\begin{warningbox}
  \textbf{Probl\`eme :} le bo\^itier \'electronique a \textbf{peu d'espace pour les
  batteries}, ce qui limite la capacit\'e embarqu\'ee et donc l'autonomie et la marge
  de tension sous charge (en lien direct avec la coupure par basse tension de la
  Section~\ref{subsec:lvc}).
\end{warningbox}

\begin{okbox}
  \textbf{Am\'elioration propos\'ee :} revoir la conception m\'ecanique du bo\^itier
  pour loger un pack de plus grande capacit\'e (ou une meilleure r\'epartition des
  cellules), en tenant compte en plus de la place qu'occuperont les nouveaux capteurs
  et les connecteurs des antennes externes. Il convient de dimensionner le pack pour
  maintenir la tension au-dessus du seuil LVC de l'ESC durant toute la mission.
\end{okbox}

% =============================================================================
\section{\'Etat actuel du syst\`eme}
% =============================================================================

\begin{longtable}{C{1.4cm} L{5.3cm} C{2.6cm} L{3.6cm}}
  \toprule
  \rowcolor{tableheadbg}
  \textcolor{tableheadfg}{\textbf{Phase}} &
  \textcolor{tableheadfg}{\textbf{Objectif}} &
  \textcolor{tableheadfg}{\textbf{\'Etat}} &
  \textcolor{tableheadfg}{\textbf{Notes}} \\
  \midrule
  \endhead
  \bottomrule
  \endfoot
  1 & Contr\^ole manuel RC & \done\ Complet & Mode unifi\'e + moteur bidirectionnel \\[3pt]
  2 & GPS + navigation & \done\ Complet & Fix test\'e en ext\'erieur \\[3pt]
  3 & LoRa + t\'el\'em\'etrie & \done\ Complet & TX/RX v\'erifi\'es \\[3pt]
  PCB V2 & Nouveau c\^ablage + I2C libre & \done\ Complet & Compil\'e et test\'e \\[3pt]
  5 & Estimation du vent & \done\ Cod\'e & Non valid\'e sur l'eau \\[3pt]
  6 & Navigation autonome & \done\ Cod\'e & Non test\'e sur l'eau, sans tests \\[3pt]
  4 & Capteurs I2C & \todo\ En attente & Voir Section~\ref{sec:abiertos} \\[3pt]
  7 & Journalisation / mission compl\`ete & \todo\ En attente & SPIFFS, propulsion auto \\
\end{longtable}

\subsection{Valeurs de calibration actuelles (r\'esum\'e)}

\begin{longtable}{L{5cm} C{2.4cm} L{5cm}}
  \toprule
  \rowcolor{tableheadbg}
  \textcolor{tableheadfg}{\textbf{Param\`etre}} &
  \textcolor{tableheadfg}{\textbf{Valeur}} &
  \textcolor{tableheadfg}{\textbf{Notes}} \\
  \midrule
  \endhead
  \bottomrule
  \endfoot
  \texttt{SAIL\_CENTER\_US} & \us{1520} & Centre Futaba (pas 1500) \\[2pt]
  \texttt{SAIL\_PLUS/MINUS\_US} & 1575 / 1465 & $\pm10^{\circ}$ --- calibrer au banc \\[2pt]
  \texttt{ROTOR\_MIN/MAX\_US} & 1417 / 1583 & $\pm90^{\circ}$ physique (manuel) \\[2pt]
  \texttt{ROTOR\_AUTO\_RANGE\_DEG} & $20^{\circ}$ & Course utilis\'ee en auto \\[2pt]
  \texttt{ESC\_REVERSE\_MIN\_US} & \us{1000} & $-100~\%$ (arri\`ere, bidirectionnel) \\[2pt]
  \texttt{ESC\_NEUTRAL\_US} & \us{1500} & Neutre / arr\^et \\[2pt]
  \texttt{ESC\_MAX\_US} & \us{2000} & $+100~\%$ (avant) \\[2pt]
  \texttt{CH3\_CENTER\_US} & \us{1545} & Centre des gaz --- v\'erifier au banc \\[2pt]
  \texttt{ESC\_SLEW\_US} & 30~\textmu s/tick & Limitation de variation \\[2pt]
  Rapport diviseur batterie & 5,683 & (562k + 120k) / 120k \\
\end{longtable}

% =============================================================================

\end{document}
